\subsection*{Defining a \emph{Spanned Data}}

Defining a \emph{spanned data} requires both the element type and its shape. For example,

\begin{lstlisting}[language=croqtile,caption={A spanned-data declaration}]
f32 [32, 16] data;
\end{lstlisting}

This defines a buffer with a shape of \texttt{[32, 16]}, where the element type is \texttt{f32}, which is the \emph{IEEE-754} single-precision float. It is also possible to use a named \texttt{mdspan} to define the buffer:

\begin{lstlisting}[language=croqtile,caption={A standalone shape}]
shape : [72, 14];
s32 [shape] data0;              // using the named mdspan `shape`
u32 [shape(0), shape(1)] data1; // same shape as the above
u16 [shape + 6] data2;          // shaped as [78, 20]
s16 [shape, 8] data3;           // shaped as [72, 14, 8]
\end{lstlisting}

This example uses the named \texttt{mdspan} \texttt{shape} to declare the two-dimensional buffer \texttt{data0} with 32-bit integer elements. Place the shape inside \texttt{[]}; element access, arithmetic, and concatenation can also appear in the shape expression.

In Croqtile terminology, the element type is referred to as the \textbf{fundamental type}. Croqtile supports multiple fundamental types, including:

\begin{itemize}[leftmargin=1.6em]
\item \textbf{Unsigned Integers}: \texttt{u8}, \texttt{u16}, \texttt{u32}
\item \textbf{Signed Integers}: \texttt{s8}, \texttt{s16}, \texttt{s32}
\item \textbf{Floating Points}: \texttt{f16}, \texttt{bf16}, \texttt{f32}
\end{itemize}

The prefixes \texttt{u} and \texttt{s} mean unsigned and signed, and the numeric suffix gives the bit width. The type \texttt{f16} is IEEE 754 half precision; \texttt{bf16} is the 16-bit brain floating-point format commonly used in machine learning.

The element type \co@s32@ and the control type \co@int@ have different roles, even though both describe 32-bit integers. A declaration that adds a shape creates a buffer; without one, the same element type declares an ordinary control integer:

\begin{lstlisting}[language=croqtile,caption={Element type without a shape}]
s32 a;    // a control integer, the same as 'int a'
s32 [4] b; // a buffer of four elements
\end{lstlisting}

Both forms compile. The shape decides which one you get, not the type name.
